\chapter{پیمایش در سیستم فایل (Navigation)}

نخستین مهارت حیاتی در محیط خط فرمان (پس از تسلط بر تایپ دستورات)، توانایی جابه‌جایی در ساختار سیستم فایل لینوکس است. در این فصل، دستورات زیر کالبدشکافی خواهند شد:

\begin{itemize}
  \item \cmd{pwd}: نمایش مسیر دایرکتوری کاری فعلی.
  \item \cmd{cd}: تغییر دایرکتوری جاری.
  \item \cmd{ls}: فهرست کردن محتویات یک دایرکتوری.
\end{itemize}

\section{ساختار درختی سیستم فایل}

لینوکس نیز همانند ویندوز، فایل‌های خود را در یک ساختار سلسله‌مراتبی (\en{Hierarchical}) سازمان‌دهی می‌کند. این بدان معناست که تمامی داده‌ها در قالب یک ساختار درختی از دایرکتوری‌ها (که در سیستم‌عامل‌های غیرفنی «پوشه» نامیده می‌شوند) چیدمان شده‌اند. هر دایرکتوری می‌تواند شامل فایل‌ها و دایرکتوری‌های فرعی (\en{Subdirectories}) دیگری باشد.

نقطه آغازین این ساختار، دایرکتوری ریشه (\en{Root Directory}) نام دارد. دایرکتوری ریشه دربرگیرنده‌ی فایل‌ها و زیرشاخه‌هایی است که خود به شاخه‌های بیشتری منشعب می‌شوند و این روند به صورت پیوسته ادامه می‌یابد.

شایان ذکر است که برخلاف ویندوز که برای هر واحد ذخیره‌سازی (مانند درایو \lr{C:} یا \lr{D:}) یک ساختار درختی مجزا تعریف می‌کند، سیستم‌های شبه‌یونیکس و لینوکس همواره دارای یک ساختار درختی واحد هستند؛ صرف‌نظر از اینکه چه تعداد هارددیسک یا دستگاه ذخیره‌سازی به رایانه متصل باشد. این تجهیزات در نقاط مختلفی از این درخت واحد «متصل» (در اصطلاح فنی \en{Mount}) می‌شوند. مدیریت این نقاط اتصال بر عهده‌ی مدیر سیستم (\en{System Administrator}) است.

\section{دایرکتوری کاری فعلی (Current Working Directory)}

اغلب کاربران با مدیر فایل گرافیکی (\en{File Manager}) تعامل دارند که ساختار درختی سیستم فایل را به‌صورت بصری نمایش می‌دهد. در این واسط‌ها، این ساختار درختی معمولاً به‌صورت وارونه ترسیم می‌شود؛ یعنی ریشه در بالاترین سطح قرار دارد و زیرشاخه‌ها به سمت پایین گسترش می‌یابند.

اما در محیط متنی خط فرمان، به دلیل نبود عناصر بصری، باید رویکرد متفاوتی برای درک این ساختار اتخاذ کرد. تصور کنید سیستم فایل مانند یک هزارتوی سه‌بعدی به شکل درخت است و شما در هر لحظه در نقطه‌ی مشخصی از این هزارتو مستقر هستید.

\bookfigure{img-01.png}{نمای درختی سیستم فایل در مدیر فایل گرافیکی}

در هر لحظه از تعامل با سیستم‌عامل، یک دایرکتوری کاری فعلی وجود دارد که به‌عنوان مبدأ فعالیت‌های شما شناخته می‌شود. این دایرکتوری شامل فایل‌ها و دایرکتوری‌های فرعی است که مستقیماً در دسترس قرار دارند. همچنین، همواره مسیری به دایرکتوری بالادستی یا والد (\en{Parent Directory}) وجود دارد. برای نمایش مسیر مطلقِ نقطه‌ای که در آن حضور دارید، از دستور \cmd{pwd} (مخفف \en{Print Working Directory}) استفاده می‌شود:

\begin{latin}
\begin{lstlisting}
[me@linuxbox ~]$ pwd
/home/me
\end{lstlisting}
\end{latin}

هنگام ورود به سیستم (\en{Login}) یا آغاز یک نشست ترمینال جدید، دایرکتوری کاری فعلی به‌طور خودکار بر روی دایرکتوری خانگی (\en{Home Directory}) شما تنظیم می‌شود. این فضا به‌طور انحصاری متعلق به حساب کاربری شماست و تنها محیطی است که یک کاربر عادی در آن مجوز تام برای ایجاد، حذف یا ویرایش فایل‌ها را دارا می‌باشد.

\section{مشاهده محتویات دایرکتوری (ls)}

برای نمایش محتویات یک دایرکتوری، اعم از فایل‌ها و دایرکتوری‌های فرعی موجود در آن، از دستور \cmd{ls} (مخفف \en{List}) استفاده می‌شود:

\begin{latin}
\begin{lstlisting}
[me@linuxbox ~]$ ls
Desktop  Documents  Music  Pictures  Public  Templates  Videos
\end{lstlisting}
\end{latin}

دستور \cmd{ls} در سیستم‌عامل‌های خانواده یونیکس، ابزاری بنیادین جهت پایش محتوای دایرکتوری‌هاست. این ابزار به کاربران اجازه می‌دهد تا محتویات هر نقطه‌ای از سیستم فایل را، فراتر از دایرکتوری کاری فعلی، مشاهده کنند. \cmd{ls} به همراه گزینه‌ها (\en{Options}) و پرچم‌های (\en{Flags}) متعدد، قابلیت‌های پیشرفته‌ای برای فیلتر کردن، قالب‌بندی خروجی و نمایش جزئیات فنی فایل‌ها (مانند مجوزها و حجم) ارائه می‌دهد که در بخش‌های آتی به‌طور جامع بررسی خواهند شد.

\section{جابه‌جایی در سلسله‌مراتب سیستم فایل (cd)}

به منظور تغییر دایرکتوری کاری فعلی و جابه‌جایی در ساختار درختی سیستم، از دستور \cmd{cd} (مخفف \en{Change Directory}) استفاده می‌شود. برای بهره‌گیری از این دستور، پس از عبارت \cmd{cd}، مسیر دایرکتوری هدف (\en{Pathname}) به عنوان آرگومان وارد می‌گردد. مسیر در واقع توالی دایرکتوری‌هایی است که باید برای رسیدن به مقصد نهایی پیموده شوند.

تعیین مسیر به دو روش کلی صورت می‌پذیرد: مسیر مطلق (\en{Absolute Path}) و مسیر نسبی (\en{Relative Path}). در ادامه، به بررسی ساختار مسیرهای مطلق خواهیم پرداخت.

\section{مسیرهای مطلق (Absolute Pathnames)}

یک مسیر مطلق از دایرکتوری ریشه (\en{Root Directory}) آغاز شده و با پیمایش سلسله‌مراتب درختی، دقیقاً به دایرکتوری یا فایل مقصد منتهی می‌شود. در لینوکس، دایرکتوری ریشه با کاراکتر اسلش (\cmd{/}) در ابتدای مسیر بازشناخته می‌شود.

به‌عنوان نمونه، در اکثر توزیع‌های لینوکس، بخش عمده‌ای از فایل‌های اجرایی و برنامه‌های سیستمی در مسیر \file{/usr/bin/} مستقر هستند. این مسیر را می‌توان به‌صورت گام‌به‌گام چنین تفسیر کرد:

\begin{enumerate}
  \item از بالاترین سطح سیستم یا همان دایرکتوری ریشه (\cmd{/}) شروع کنید.
  \item وارد زیرشاخه‌ی \file{usr} شوید.
  \item در نهایت به دایرکتوری \file{bin} که درون \file{usr} قرار دارد، برسید.
\end{enumerate}

به‌طور کلی، هر مسیری که با کاراکتر \cmd{/} آغاز شود، یک مسیر مطلق تلقی می‌گردد؛ زیرا موقعیت هدف را به‌صورت قطعی و مستقل از مکان فعلی کاربر، نسبت به ریشه‌ی سیستم تعیین می‌کند.

\begin{latin}
\begin{lstlisting}
[me@linuxbox ~]$ cd /usr/bin
[me@linuxbox bin]$ pwd
/usr/bin
[me@linuxbox bin]$ ls
...Listing of many, many files ...
\end{lstlisting}
\end{latin}

همان‌طور که ملاحظه می‌کنید، دایرکتوری کاری فعلی به \file{/usr/bin/} تغییر یافته و محتویات آن فهرست شده است. اعلان خط فرمان (\en{Prompt}) نیز این جابه‌جایی را منعکس کرده است. نمایش موقعیت فعلی در اعلان، یکی از ویژگی‌های کلیدی و کاربردی در محیط‌های پوسته (\en{Shell}) است که به کاربر کمک می‌کند همواره از جایگاه خود در ساختار پیچیده‌ی سیستم فایل آگاه باشد.

\section{مسیرهای نسبی (Relative Pathnames)}

برخلاف مسیرهای مطلق که همواره از دایرکتوری ریشه (\en{Root}) آغاز می‌شوند، مسیرهای نسبی مبدأ خود را دایرکتوری کاری فعلی (\en{Current Working Directory}) قرار می‌دهند. برای پیمایش نسبی در ساختار سیستم فایل، از دو نشانگر قراردادی و بنیادین استفاده می‌شود:

\begin{itemize}
  \item نقطه (\cmd{.}): به دایرکتوری کاری فعلی اشاره دارد.
  \item دو نقطه (\cmd{..}): به دایرکتوری والد (\en{Parent Directory}) اشاره می‌کند.
\end{itemize}

برای درک بهتر عملکرد این نشانگرها، ابتدا دایرکتوری کاری را به \file{/usr/bin/} تغییر می‌دهیم:

\begin{latin}
\begin{lstlisting}
[me@linuxbox ~]$ cd /usr/bin
[me@linuxbox bin]$ pwd
/usr/bin
\end{lstlisting}
\end{latin}

با فرض اینکه در دایرکتوری \file{/usr/bin/} مستقر باشیم، برای بازگشت به دایرکتوری والد (یعنی \file{/usr/})، دو روش پیش رو داریم:

\begin{enumerate}
  \item استفاده از مسیر مطلق:
\begin{latin}
\begin{lstlisting}
[me@linuxbox bin]$ cd /usr
[me@linuxbox usr]$ pwd
/usr
\end{lstlisting}
\end{latin}

  \item استفاده از مسیر نسبی (بهینه):
\begin{latin}
\begin{lstlisting}
[me@linuxbox bin]$ cd ..
[me@linuxbox usr]$ pwd
/usr
\end{lstlisting}
\end{latin}
\end{enumerate}

در دنیای لینوکس، همواره روشی که مستلزم تایپ کاراکترهای کمتر است، کارآمدتر و حرفه‌ای‌تر تلقی می‌شود. به‌عنوان مثال، برای بازگشت به دایرکتوری \file{bin} از مبدأ \file{/usr/} می‌توانیم از هر دو شیوه استفاده کنیم.

\begin{enumerate}
  \item از طریق مسیر مطلق:
\begin{latin}
\begin{lstlisting}
[me@linuxbox usr]$ cd /usr/bin
[me@linuxbox bin]$ pwd
/usr/bin
\end{lstlisting}
\end{latin}

  \item یا از طریق مسیر نسبی:
\begin{latin}
\begin{lstlisting}
[me@linuxbox usr]$ cd ./bin
[me@linuxbox bin]$ pwd
/usr/bin
\end{lstlisting}
\end{latin}
\end{enumerate}

اکنون به یک قاعده مهم در ترمینولوژی پوسته می‌پردازیم: در اکثر سناریوها، نشانگر نقطه (\cmd{.}) که به دایرکتوری جاری اشاره دارد، به‌صورت ضمنی در نظر گرفته شده و قابل حذف است. این بدان معناست که برای ورود به زیرشاخه‌ای که در دایرکتوری فعلی قرار دارد، نیازی به ذکر \file{./} نیست:

\begin{latin}
\begin{lstlisting}
[me@linuxbox usr]$ cd bin
\end{lstlisting}
\end{latin}

بنابراین، دو دستور \cmd{cd ./bin} و \cmd{cd bin} از نظر عملکردی کاملاً یکسان هستند، اما دستور دوم به‌دلیل اختصار، گزینه‌ی ارجح برای کاربران حرفه‌ای است.

به منظور تسریع در پیمایش سلسله‌مراتب سیستم فایل، دستور \cmd{cd} از میان‌برهای قراردادی بهره می‌برد که امکان جابه‌جایی آنی به مسیرهای کلیدی را فراهم می‌کنند. در جدول زیر، این پارامترها و عملکرد عملیاتی هر یک تشریح شده است:

\begin{table}[htbp]
\centering
\caption{میان‌برهای کاربردی دستور \cmd{cd}}
\label{tab:cd-shortcuts}
\begin{tabularx}{\textwidth}{l X}
\toprule
\textbf{میان‌بر} & \textbf{نتیجه عملیاتی} \\
\midrule
\cmd{cd} & دایرکتوری کاری را مستقیماً به دایرکتوری خانگی (\en{Home Directory}) کاربر جاری تغییر می‌دهد. \\
\addlinespace
\cmd{cd -} & دایرکتوری کاری را به موقعیت قبلی بازمی‌گرداند؛ این میان‌بر برای جابه‌جایی سریع میان دو مسیر متفاوت بسیار کارآمد است. \\
\addlinespace
\cmd{cd \textasciitilde user\_name} & مسیر را به دایرکتوری خانگی کاربر مشخص‌شده تغییر می‌دهد. برای مثال، \cmd{cd \textasciitilde bob} شما را به دایرکتوری خانگی کاربر «\en{bob}» منتقل می‌کند. \\
\bottomrule
\end{tabularx}
\end{table}

\section{اصول و ضوابط نام‌گذاری فایل‌ها و دایرکتوری‌ها}

نام‌گذاری فایل‌ها و دایرکتوری‌ها در لینوکس، با وجود شباهت به سایر سیستم‌عامل‌ها، تابع منطق منحصربه‌فردی است که درک آن برای مدیریت حرفه‌ای سیستم ضرورت دارد:

\begin{enumerate}
  \item \textbf{فایل‌های پنهان (\en{Hidden Files}):} فایل‌هایی که نام آن‌ها با کاراکتر نقطه (\cmd{.}) آغاز می‌شود، به‌صورت پیش‌فرض در خروجی استاندارد دستور \cmd{ls} ظاهر نمی‌شوند. این فایل‌ها که عمدتاً حاوی تنظیمات نرم‌افزاری و پیکربندی‌های پروفایل کاربری هستند، تنها با استفاده از گزینه \cmd{-a} (مانند \cmd{ls -a}) قابل مشاهده‌اند. در بدو ایجاد هر حساب کاربری، مجموعه‌ای از این فایل‌ها جهت تنظیم محیط پوسته به‌صورت خودکار در دایرکتوری خانگی مستقر می‌شوند.

  \item \textbf{حساسیت به حروف:} برخلاف محیط‌هایی نظیر ویندوز، لینوکس میان حروف بزرگ و کوچک تمایز قائل می‌شود. از این منظر، \file{File1} و \file{file1} دو هویت مستقل و مجزا در سطح سیستم فایل تلقی می‌شوند. این ویژگی، دقت بالایی را در فراخوانی دستورات و مدیریت داده‌ها ایجاب می‌کند.

  \item \textbf{استقلال از پسوندهای فایل:} در معماری لینوکس، مفهوم «پسوند فایل» (\en{File Extension}) به معنای فنی و اجباری آن وجود ندارد. ماهیت و نوع یک فایل معمولاً از طریق فراداده‌ها یا اعداد جادویی (\en{Magic Numbers}) شناسایی می‌شود. با این حال، استفاده از پسوندهای قراردادی (مانند \file{.jpg} یا \file{.txt}) صرفاً جهت تسهیل در شناسایی بصری توسط کاربر و تعامل بهتر با برخی برنامه‌های گرافیکی توصیه می‌گردد.

  \item \textbf{استانداردهای پیشنهادی نام‌گذاری:} لینوکس اجازه استفاده از نام‌های طولانی و کاراکترهای متنوع را صادر می‌کند، اما جهت پیشگیری از تداخلات فنی، به‌ویژه در فرآیندهای خودکارسازی و اسکریپت‌نویسی، رعایت موارد زیر الزامی است:
  \begin{itemize}
    \item از درج فاصله (\en{Space}) در نام فایل‌ها جداً خودداری کنید؛ به‌جای آن از کاراکتر زیرخط (\cmd{\_}) استفاده نمایید.
    \item استفاده از کاراکترهای خاص را محدود کرده و برای جداسازی صرفاً به نقطه (\cmd{.}) و خط تیره (\cmd{-}) بسنده کنید.
  \end{itemize}
\end{enumerate}

رعایت این رویکرد انضباطی، ضامن ثبات سیستم و جلوگیری از بروز رفتارهای پیش‌بینی‌نشده در هنگام پردازش‌های دسته‌ای است.

\section{جمع‌بندی}

در این فصل، به بررسی ساختار درختی دایرکتوری‌ها در سیستم‌های لینوکسی و نحوه‌ی تعامل پوسته (\en{Shell}) با این معماری پرداختیم. مفاهیم کلیدی از جمله مسیرهای مطلق (\en{Absolute Paths}) و مسیرهای نسبی (\en{Relative Paths}) تبیین و تفاوت‌های بنیادین آن‌ها در فرآیند آدرس‌دهی تشریح شد. افزون بر این، با تسلط بر دستور \cmd{cd} و میان‌برهای عملیاتی آن، رویکردهای کارآمد جهت پیمایش سریع در این ساختار سلسله‌مراتبی را فراگرفتیم.

با کسب این دانش زیربنایی، در فصل آتی به کاوش در اجزای یک سیستم لینوکس مدرن خواهیم پرداخت و از مهارت‌های آموخته‌شده برای پیمایش عمیق‌تر و مدیریت بهینه‌ی فایل‌ها بهره خواهیم جست.